% DCCA-GS -- ICRA 2027 submission (IEEE conference template).
% Non-experiment parts (abstract / introduction / related work / method) are complete;
% the experiments section is a skeleton pending the experiment matrix (see
% docs/05-paper/后续实验计划.md and the HAC++-replication plan).
% ICRA page budget: 6 pages of content + 2 allowed with fee; references extra.
\documentclass[conference]{IEEEtran}
\IEEEoverridecommandlockouts
\usepackage{cite}
\usepackage{amsmath,amssymb,amsfonts}
\usepackage{algorithmic}
\usepackage{algorithm}
\usepackage{graphicx}
\usepackage{textcomp}
\usepackage[dvipsnames]{xcolor}
\usepackage{booktabs}
\usepackage{multirow}
\usepackage{colortbl}
\usepackage[normalem]{ulem}
\usepackage{balance}

\newcommand{\DCCA}{DCCA-GS}
\newcommand{\todo}[1]{\textcolor{red}{[TODO: #1]}}

\begin{document}

\title{DCCA-GS: Decoder-Reproducible Content-Adaptive Quantization and\\
Coverage-Aware Fusion Pruning for Anchor-Based 3D Gaussian Splatting}

\author{\IEEEauthorblockN{Chen Tong}
\IEEEauthorblockA{Department of ..., University of ...\\
Email: todo@example.edu}}

\maketitle

\begin{abstract}
Anchor-based codecs such as HAC++ compress 3D Gaussian Splatting (3DGS) scenes by
more than an order of magnitude, yet two decisions inside them remain crude. First,
the quantization step of each anchor is predicted from hash context alone: the
allocation ignores how complex the local content is and how much it matters to the
rendered image, and carrying per-anchor step decisions in the bitstream would violate
the zero-side-information contract that makes anchor codecs practical. Second,
anchor pruning selects survivors with static, per-anchor scores, which is blind to
spatial coverage: under a tight budget the score keeps high-sensitivity detail
anchors and deletes the few anchors that cover large background regions, and
reconstruction collapses. We present \DCCA{}, which addresses both problems with two
decoder-reproducible mechanisms. (1)~A \emph{complexity multiplier} rescales the
quantization step of every anchor by a factor computed from four structural
statistics that the decoder can recompute exactly; during training the same small
network is supervised by render-sensitivity statistics, so fine steps flow to anchors
that matter to the image---without a single bit of side information. (2)~\emph{coverage-aware
fusion pruning} augments the ADMM sparsification of GaussianSpa with a
coverage-dominant score that fuses per-anchor coverage with render sensitivity, and
with a hard projection that keeps at least one anchor per occupied coarse cell. On
three large-scale UAV scenes, \DCCA{} needs on average $58\%$ of the HAC++ model size
at matched $\lambda$ while remaining within $0.27$--$0.61$~dB PSNR, ranks first
among 18 compared methods on a per-scene rank aggregation, and its coverage
guards eliminate the failure mode behind a $1.56$~dB low-rate collapse at a
quality cost within measurement noise.
\end{abstract}

\section{Introduction}

3D Gaussian Splatting (3DGS)~\cite{kerbl2023} represents a scene with explicit
Gaussian primitives and renders novel views in real time, but high-quality scenes
require millions of primitives whose attributes dominate storage. Anchor-based
codecs reduce this cost by organizing Gaussians around a compact set of anchors:
Scaffold-GS~\cite{lu2024} derives neural Gaussians from anchor features,
HAC~\cite{chen2024} and HAC++~\cite{chen2025} add hash-grid contexts and
conditional entropy models, and ContextGS~\cite{wang2024b} models anchor-level
context. These systems push the rate--distortion frontier far beyond vanilla 3DGS,
and we build directly on the HAC++ coding pipeline. Two mechanisms inside this
pipeline, however, allocate the remaining bits crudely, and both are the subject of
this paper.

\subsection{Bit allocation is content-blind}

HAC++ predicts a \emph{base} quantization step per anchor from hash-grid context
(adaptive quantization module, AQM). The base step, however, is modulated neither
by the complexity of the local content nor by the importance of the anchor to the
rendered image: a flat wall and a cluttered facade that share the same hash context
receive the same step. Making the step content-adaptive seems to require either an
explicit per-anchor step array---side information that changes the coding contract
and costs bits---or latents whose rate must itself be modeled. Our setting forbids
both: we require \emph{zero side information}, i.e., every signal that influences
quantization or rate estimation must be recomputable by the decoder from the decoded
bitstream alone.

The way out is to separate \emph{what the step depends on} from \emph{what guides
it}. Structural statistics---local anchor density, predicted scaling anisotropy,
predicted offset energy, and active mask ratio---are functions of quantities the
decoder already has, so a small network over these statistics is
\emph{decoder-reproducible} by construction. What this network should \emph{predict}
is not directly observable at test time, but during training we can measure how much
each anchor matters: the gradient of the rendering loss with respect to its
attributes. We therefore train the shared network to regress a render-sensitivity
target accumulated as an exponential moving average. The supervision exists only in
the training objective; the bitstream and the decode path are untouched, and the
encoder--decoder round trip remains bit-exact.

\subsection{Pruning is coverage-blind}

The coarsest rate lever in an anchor codec is the number of surviving anchors.
Training-side sparsification~\cite{zhang2025gaussianspa} imposes a hard budget
during optimization via ADMM, so survivors re-adapt to the final budget---strictly
better than pruning after training. Yet the selection step itself is still a greedy
top-$k$ over \emph{static, per-anchor} scores. Such scores cannot express the
statement ``some anchor must remain in this region'': under a tight budget, a score
dominated by render sensitivity retains cluttered detail anchors and removes the
sparse anchors that cover large background areas. The result is a \emph{coverage
hole}: in our experiments, fusing sensitivity into the pruning score without any
coverage term collapses low-rate reconstruction by $1.56$~dB on scene 1-78 while
leaving the anchor count almost unchanged---the pruner did not prune too much, it
pruned the wrong anchors.

We fix this structurally rather than by re-tuning weights. The fused pruning score
is made \emph{coverage-dominant} (sensitivity is capped at a $0.3$ share), and the
ADMM hard projection is extended with a coverage constraint: every occupied coarse
cell of the scene keeps at least one anchor. The constraint is a one-line change to
the projection step, costs nothing at decode time, and eliminates the collapse
failure mode. A pre-registered re-validation of the repaired score improves the
tightest-budget point and holds size parity, but does not clear our improvement
gates; we therefore report coverage-aware fusion pruning as a robustness
mechanism rather than a rate--distortion gain, and analyze the signal defects
behind this boundary in Section~\ref{sec:analysis}.

\subsection{Contributions}

\begin{itemize}
\item \textbf{Decoder-reproducible content-adaptive quantization.} A complexity
multiplier $m_i$ rescales the AQM step of anchor $i$ from four structural
statistics; a single shared $8{\to}32{\to}3$ network produces it, is supervised in
training by render-sensitivity gradients, and is recomputed bit-exactly by the
decoder. The bitstream carries no step array and no side information. Removing the
multiplier costs $0.41$~dB at $\lambda{=}0.002$ on 1-78.
\item \textbf{Coverage-aware fusion pruning.} On top of the ADMM budget
sparsification of GaussianSpa~\cite{zhang2025gaussianspa}, which we adopt as-is as
the base mechanism, we add (i) depth re-initialization at growth stop, which inserts
surface candidates that must compete under the \emph{existing} budget, and (ii) a
sensitivity-aware pruning strategy: a coverage-dominant fusion score plus an ADMM
coverage constraint. The strategy removes the failure mode behind the measured
$1.56$~dB low-rate collapse at a quality cost within measurement noise;
Section~\ref{sec:analysis} quantifies this boundary and dissects the signal
defects that limit the rate--distortion gains of importance-guided selection.
\item We validate both mechanisms on three large-scale UAV scenes under a uniform
protocol ($30$k iterations, seed $42$, fp32 sizes, bit-exact decode), and document
the negative results that shaped the design.
\end{itemize}

\begin{figure}[t]
\centering
\todo{overview figure: left = anchor codec pipeline with the complexity multiplier
hooked into AQM; right = SPA projection step with the coverage constraint and the
fused score}
\caption{\DCCA{} overview. Both innovations are decoder-reproducible: the
complexity multiplier lives in the quantization path, the coverage-aware fusion
pruning lives in the training-side projection. Neither writes anything extra to the
bitstream.}
\label{fig:overview}
\end{figure}

\section{Related Work}

\subsection{Compression of 3D Gaussian Splatting}
Compact-3DGS methods reduce the primitive count and the per-primitive cost:
LightGaussian~\cite{fan2024} distills and prunes by global significance,
Compact3DGS~\cite{lee2024} learns a compact mask, Reduced3DGS~\cite{papantonakis2024}
sensitively quantizes attributes, MesonGS~\cite{xie2024} transforms attributes
post-training, and PCGS~\cite{chen2026a} refines masks and quantization
progressively. Anchor-based codecs organize primitives around shared anchors:
Scaffold-GS~\cite{lu2024} introduces view-adaptive anchor features,
HAC~\cite{chen2024} couples anchors with hash-grid contexts, HAC++~\cite{chen2025}
adds intra-anchor context with a Gaussian mixture likelihood, and
ContextGS~\cite{wang2024b} models anchor-level context autoregressively. \DCCA{} is
implemented on the open-source gsplat rasterization library~\cite{ye2024gsplat} and
reproduces the HAC++ coding pipeline within it; all HAC++ reference numbers used for
comparison are produced by the official HAC++ implementation under an identical
protocol. \DCCA{} is
orthogonal to these probability-modeling improvements: it changes \emph{which} step
each anchor is quantized with and \emph{which} anchors survive, while keeping the
HAC++ bitstream contract.

\subsection{Rate--Distortion Optimization and Adaptive Quantization}
Learned image compression optimizes a differentiable rate estimate jointly with
distortion~\cite{balle2018,cheng2020} and allocates precision through learned
latents. RDO-Gaussian~\cite{wang2024a}, CompGS~\cite{liu2024b}, and
CAT-3DGS~\cite{zhan2025} transfer this recipe to 3DGS. In all of them, per-element
precision is carried by the latent itself and its cost by the entropy model. Our
constraint is stricter: no signal beyond the decoded bitstream may influence
quantization. The complexity multiplier therefore consumes only structural
statistics and is guided, in training only, by render sensitivity---a mechanism
reminiscent of perceptual JND (just-noticeable-difference) weighting in image
codecs, but realized as a decoder-reproducible network rather than a side
information channel.

\subsection{Sparsification and Anchor Placement}
HAC~\cite{chen2024} and HAC++~\cite{chen2025} promote offset sparsity with mask
losses folded into rate estimation. GaussianSpa~\cite{zhang2025gaussianspa}
sparsifies during training by alternating a smooth sub-problem with a hard top-$k$
projection---the ADMM structure we adopt as the budget mechanism. Placement has
been studied by Mini-Splatting~\cite{fang2024}, which re-initializes Gaussians from
depth back-projection and simplifies by contribution; we borrow only its
depth-reinitialization idea, run it once at growth stop, and force the new
candidates to compete under the existing ADMM budget instead of growing the model.
What is missing in all of these selection rules is spatial coverage: scores are
per-anchor and static, so no rule can express ``this region must stay covered''.
Our fusion pruning contributes exactly this missing piece, both softly
(coverage-dominant score) and structurally (a hard projection constraint).

\section{Method}

\subsection{Preliminaries and Coding Contract}
We build on an anchor-based codec with $N$ anchors and $K$ offsets per anchor. Each
anchor carries features $a_i\in\mathbb{R}^{D_a}$, scaling parameters, offsets, and a
binary offset mask; each valid offset instantiates one neural Gaussian whose color
and opacity are decoded by shared MLPs from a hash grid. The codec is trained with
the Lagrangian objective
\begin{equation}
\min_{\theta}\; D(\theta) + \lambda\, R(\theta),
\label{eq:rd}
\end{equation}
where $\theta$ collects all coded parameters and $R$ is a differentiable rate
estimate. Following HAC++, the base quantization step of anchor $i$ for field $f$
(feature, scaling, or offset) is predicted from hash context by the AQM,
$Q^{\mathrm{AQM}}_{i,f}$.

\emph{Zero-side-information contract.} The decoder may use: (i) the decoded
bitstream, (ii) the stored entropy-model and decoder weights, and (iii) any function
of already-decoded quantities (hash grid, anchor geometry, masks, decoded MLP
outputs). Anything else---per-anchor steps, importance maps, sensitivity
maps---is forbidden. Both mechanisms below respect this contract by construction;
the encoder--decoder round trip is verified bit-exactly for every reported run.

\subsection{Complexity Multiplier with Render-Sensitivity Supervision}
\label{sec:multiplier}

\subsubsection{Multiplier}
The final step used to quantize anchor $i$, field $f$, is
\begin{equation}
Q_{i,f} \;=\; Q^{\mathrm{AQM}}_{i,f}\,
\bigl(1+\tanh q^{\mathrm{AQM}}_{i,f}\bigr)\,
\bigl(1+\alpha\,\tanh m_{i,f}\bigr),
\label{eq:q}
\end{equation}
where $\alpha$ ramps from $0$ to a configured scale during training so that
quantization-aware training starts from the HAC++ behaviour. The multiplier is
produced by a small shared network, $m_{i,f}=[\mathrm{MLP}_c(x_i)]_{f}$ with
$\mathrm{MLP}_c: 8{\to}32{\to}3$, from the input vector $x_i$ of four structural
statistics: \emph{local anchor density} $\rho_i$ (occupancy of the anchor's voxel
neighbourhood), \emph{predicted scaling anisotropy} (spread of the decoded
$\log$-scales), \emph{predicted offset energy} (norm of the decoded offsets), and
\emph{active mask ratio} (fraction of surviving offsets). All four are functions of
quantities the decoder possesses after decoding the grid and the anchor attributes,
so $\mathrm{MLP}_c$---whose weights are stored in the bitstream as part of the
entropy model---is evaluated identically at both ends. This is the same property
that makes AQM itself decoder-reproducible; we extend it from a context-derived
\emph{base} step to a content-adaptive \emph{modulation} of that step. Intuitively,
$\mathrm{MLP}_c$ lowers the step (finer quantization) where content is dense,
anisotropic, and high-energy, and raises it where few, isotropic, low-energy
anchors cover smooth regions.

\subsubsection{Render-sensitivity supervision}
Nothing in $x_i$ says whether the anchor \emph{matters}. During training we measure
exactly that: retaining gradients through the quantization-aware forward pass, we
accumulate per-anchor sensitivity
\begin{equation}
s_i \;=\; \mathrm{EMA}_\tau\Bigl(\bigl\lVert \partial \mathcal{L}_{\mathrm{render}}
/ \partial u_i \bigr\rVert_1\Bigr),
\label{eq:sens}
\end{equation}
where $u_i$ collects the pre-quantization attributes generated by anchor $i$ and
$\tau$ is the EMA decay. A bounded, per-frame-normalized target is derived from
$s_i$, and $\mathrm{MLP}_c$ is trained with an additional MSE regression term
weighted by $\lambda_{\mathrm{sens}}$. The target exists only in the training
objective: at test time the decoder runs $\mathrm{MLP}_c$ on $x_i$ alone. The
network therefore learns to \emph{predict} sensitivity from structure, and the
prediction modulates the step---fine steps flow to anchors that matter to the
rendered image, coarse steps to anchors that do not.

\subsection{Coverage-Aware Fusion Pruning}
\label{sec:fusion}

\subsubsection{Base mechanism: ADMM budget sparsification}
We adopt the training-side sparsification of GaussianSpa~\cite{zhang2025gaussianspa}
as-is. With $a_i$ the anchor mask score, a binary state $z_i$, dual variable $u_i$,
and budget $\kappa$, each projection step solves
\begin{equation}
z=\mathrm{TopK}(a+u,\ \kappa),\qquad
u=\mathrm{clamp}\!\left(u+a-z,\ \pm 1\right),
\label{eq:spa}
\end{equation}
and anchors with $z_i=0$ are pruned while the survivors continue to train under the
budget. The budget is ramped down from the historical maximum anchor count during
the growth phase. This mechanism decides \emph{how many} anchors survive; our
additions decide \emph{which} ones.

\subsubsection{Depth re-initialization under the fixed budget}
Anchor growth stops early in training, and the anchors are clustered where initial
SfM points happen to be. Following the depth-reinitialization idea of
Mini-Splatting~\cite{fang2024}---and \emph{only} that idea---we render depth maps
for $V$ sampled training cameras at growth stop, back-project them to surface
points, voxelize them, and append the highest-coverage candidates as new anchors.
Crucially, the SPA budget $\kappa$ is \emph{not} increased: the new candidates must
displace existing anchors through the projection in \eqref{eq:spa}. This lets the
anchor set migrate toward the true scene surface without growing the model.

\subsubsection{Sensitivity-aware fusion score}
The projection in \eqref{eq:spa} ranks anchors by the static mask score $a_i$.
We replace the ranking by a fused importance
\begin{equation}
\hat{s}_i \;=\; (1-w)\,c_i \;+\; w\,t_i \;+\; \beta\, a_i ,
\qquad w \le w_{\max}=0.3,
\label{eq:fusion}
\end{equation}
where $c_i$ is the \emph{coverage} of anchor $i$ (its accumulated projected area in
the rasterizer's per-anchor bookkeeping, normalized), $t_i$ is the
\emph{render sensitivity} of Section~\ref{sec:multiplier} (normalized), and $\beta$
keeps the original mask score as a tiebreaker. The cap $w_{\max}$ is the load-bearing
detail: an unconstrained $w$ lets sensitivity dominate, and we measured the
consequence---on 1-78 at $\lambda{=}0.004$, an $w{=}0.5$ fusion collapses decoded
PSNR from $27.51$ to $25.95$~dB at an almost unchanged anchor count
($507.8$k vs.\ $506.0$k): the pruner kept high-sensitivity detail anchors and
deleted low-sensitivity anchors that covered large background regions, opening
coverage holes. With $w\le0.3$, coverage is always the dominant term.

\subsubsection{ADMM coverage constraint}
The soft cap is complemented by a structural guarantee. Partition the scene into
coarse cells of size $\delta$ (we use $\delta=0.05$ in scene units); at each hard
projection, every occupied cell must keep at least one anchor:
\begin{equation}
\sum_{i\in C} z_i \;\ge\; 1 \quad \forall\, C \text{ occupied},
\label{eq:cov}
\end{equation}
implemented by protecting the best-scoring anchor of each occupied cell before the
top-$k$ of \eqref{eq:spa}. The constraint costs nothing at decode time and is
unconditional: whatever the fused scores say, no region of the scene can go
unrepresented. In our measurements the constraint alone rarely changes the anchor
count (e.g., $671{,}570$ vs.\ $671{,}497$ on 4-10); its value is that the failure
mode of \eqref{eq:fusion} is impossible by construction, which is what allowed the
fusion to be enabled at the low-rate operating points at all.

\begin{algorithm}[t]
\caption{\DCCA{} training loop (per iteration)}
\label{alg:loop}
\begin{algorithmic}[1]
\STATE render, compute $\mathcal{L}=\mathcal{L}_{\mathrm{render}}+\lambda R + \lambda_{\mathrm{sens}}\lVert \mathrm{MLP}_c(x)-s\rVert^2$ \COMMENT{Eq.~\eqref{eq:sens}}
\STATE update anchor growth; accumulate coverage $c_i$ and sensitivity $s_i$
\IF{ADMM projection step}
  \STATE $\hat{s}_i=(1-w)c_i + w t_i + \beta a_i$ \COMMENT{Eq.~\eqref{eq:fusion}, $w\le0.3$}
  \STATE protect $\arg\max_{i\in C}\hat{s}_i$ per occupied cell $C$; project $z=\mathrm{TopK}(\hat{s}+u,\kappa)$ \COMMENT{Eq.~\eqref{eq:cov}}
  \STATE update $u$; prune $z_i=0$
\ENDIF
\IF{growth stop $T_r$}
  \STATE back-project sampled depths; append voxelized candidates under fixed $\kappa$
\ENDIF
\STATE quantize with $Q_{i,f}$ \COMMENT{Eq.~\eqref{eq:q}; $\mathrm{MLP}_c$ shared, decoder-recomputed}
\end{algorithmic}
\end{algorithm}

\subsection{What the Decoder Recomputes}
Table~\ref{tab:contract} audits the coding contract. Both innovations consume only
decoder-recomputable signals at test time; the sensitivity target and the coverage
statistics are training-side only and are never encoded. The complete encode
$\to$ decode $\to$ evaluate pipeline is exercised on every reported operating
point, and a bit-exact round-trip diagnostic confirms that the decoded model is
identical to the one the encoder measured.

\begin{table}[t]
\centering
\caption{Zero-side-information audit.}
\label{tab:contract}
\begin{tabular}{p{0.31\linewidth}p{0.28\linewidth}p{0.22\linewidth}}
\toprule
Signal & Producer & In bitstream? \\
\midrule
Base step $Q^{\mathrm{AQM}}$ & hash context (HAC++) & weights only \\
Complexity input $x_i$ & decoded grid/anchors/masks & no \\
Multiplier $m_{i,f}$ & shared $\mathrm{MLP}_c$ & weights only \\
Sensitivity target $s_i$ & training gradients & training only \\
Coverage $c_i$ & rasterizer bookkeeping & training only \\
Fusion score $\hat{s}_i$ & training-side projection & training only \\
\bottomrule
\end{tabular}
\end{table}

\section{Experiments}
\label{sec:experiments}

\subsection{Setup}
We evaluate on three large-scale UAV scenes (1-78, 2-06, 4-10; 1{,}200 images
each, 150 validation views). The default protocol trains for 30k iterations
(60k in the iteration study) with feature dimension 32, 10 offsets, seed 42,
$\mathrm{data\_factor}{=}1$, 1600-px width, and one validation pass out of
every 8 views. All sizes are decode-required payload with fp32 decoder weights
(HAC++ accounting), and every operating point is produced by the full
encode\,$\to$\,decode\,$\to$\,evaluate pipeline with a bit-exact round-trip
diagnostic. Comparison numbers come from two sources: (i)~a shared evaluation
table covering 18 methods on the same scenes\,—\,compiled by a third party;
the per-method protocol is not re-verified, so we treat it as a breadth
reference and flag it accordingly; and (ii)~our own runs of the
\emph{official} HAC++ implementation under an identical protocol on scene
1-78, which we use as the rigorous reference. \DCCA{}-base denotes the
complexity multiplier (I2+I6) on top of SPA and depth re-initialization;
\DCCA{}-full additionally enables coverage-aware fusion pruning.

\subsection{Main comparison}
Table~\ref{tab:main} reports one operating point per method (PSNR / size in
MB). On 1-78, the \DCCA{}-base point dominates every entry of the shared
table on both axes simultaneously; on 2-06 and 4-10 it is the smallest model
within 0.42 and 0.06~dB of the best entry, respectively. Aggregating the
three scenes with the rank score
$\mathrm{rank}(\mathrm{PSNR})/6+\mathrm{rank}(\mathrm{SSIM})/6+
\mathrm{rank}(\mathrm{LPIPS})/6+\mathrm{rank}(\mathrm{Size})/2$ (rank~1
best, averaged over scenes), \DCCA{} scores 3.17 versus 3.56 for HAC++, the
best of the 18 compared methods. Against the stricter, same-protocol
official HAC++ runs on 1-78 (27.84~dB at 20.0~MB, 28.21 at 26.3, 28.68 at
40.6), \DCCA{}-base at 30k trades quality for size
($-0.27$ to $-0.61$~dB at $58\%$ of the size at matched $\lambda$).

\begin{table*}[t]
\centering
\caption{Decoded-view PSNR (dB) / model size (MB) on three large-scale UAV
scenes, one operating point per method. Shared-table entries are single
points compiled by a third-party evaluation (protocol not re-verified);
\DCCA{} rows are decoded bitstreams under the uniform 30k/seed-42/fp32
protocol at $\lambda{=}0.002$.}
\label{tab:main}
\small
\begin{tabular}{lccc}
\toprule
Method & 1-78 & 2-06 & 4-10 \\
\midrule
HAC++~\cite{chen2025} & 27.14 / 17.2 & 28.51 / 21.0 & 28.76 / 17.4 \\
PC-GS & 23.09 / 12.0 & 26.26 / 15.1 & 26.36 / 11.8 \\
CAT & 23.40 / 14.3 & 25.72 / 15.4 & 26.29 / 13.1 \\
ContextGS~\cite{wang2024b} & 24.90 / 29.3 & 27.20 / 26.8 & 27.39 / 22.6 \\
HAC~\cite{chen2024} & 24.83 / 21.7 & 26.86 / 26.8 & 26.94 / 20.9 \\
RDO-Gaussian~\cite{wang2024a} & 23.91 / 11.2 & 25.14 / 14.7 & 24.68 / 6.4 \\
Scaffold-GS~\cite{lu2024} & 24.63 / 286.1 & 27.21 / 250.2 & 27.21 / 202.0 \\
Compact3D & 24.27 / 21.5 & 25.68 / 33.6 & 25.53 / 19.9 \\
gsplat~\cite{ye2024gsplat} (uncompressed) & 26.02 / 604 & 27.63 / 606 & 26.82 / 610 \\
3DGS~\cite{kerbl2023} (uncompressed) & 26.06 / 604 & 27.76 / 606 & 27.02 / 610 \\
TC-GS & 23.41 / 17.7 & 24.84 / 15.7 & 25.99 / 18.6 \\
Compact3DGS~\cite{lee2024} & 24.38 / 59.6 & 25.75 / 86.1 & 25.40 / 51.4 \\
Mini-Splatting~\cite{fang2024} & 23.60 / 21.3 & 25.11 / 30.3 & 25.20 / 28.6 \\
FC-GS & 19.51 / 47.1 & 20.25 / 59.6 & 21.94 / 59.7 \\
LightGaussian~\cite{fan2024} & 24.35 / 281.7 & 26.22 / 407.5 & 25.46 / 194.9 \\
GaussianSpa~\cite{zhang2025gaussianspa} & 24.37 / 106.5 & 25.13 / 107.2 & 25.50 / 99.6 \\
Octree-GS & 20.33 / 21.5 & 21.29 / 15.7 & 24.55 / 34.7 \\
\midrule
\DCCA{}-base (I2+I6) & 27.75 / 15.0 & 28.09 / 16.4 & 28.70 / 15.7 \\
\bottomrule
\end{tabular}
\end{table*}

\subsection{Training iterations}
Doubling the budget from 30k to 60k improves \DCCA{}-base by $+0.44$ dB at
$\lambda{=}0.004$ (27.95~dB at 12.5~MB), $+0.62$~dB at $\lambda{=}0.002$
(28.37 at 16.7), and $+0.23$~dB at $\lambda{=}0.0005$ (28.18 at 25.8).
At 60k the $\lambda{=}0.004$ and $\lambda{=}0.002$ points dominate the
\emph{30k} HAC++ references on both axes; we emphasize the iteration mismatch:
HAC++ likewise improves substantially with 60k
(their reported BD-rate gain is $-32.9\%$), so this is a budget-matched
trend, not a claim against HAC++ at 60k.

\subsection{Analysis: importance-guided pruning under sparse coverage}
\label{sec:analysis}
With fusion pruning engaged (verified by per-projection engagement logging in
$>89\%$ of projections), a six-point scoreboard against \DCCA{}-base shows
the mechanism at parity overall ($-0.20$, $+0.10$, $+0.08$~dB at 30k for
$\lambda{=}0.004/0.002/0.0005$; $-0.01$, $-0.23$, $-0.12$ at 60k), with a
consistent pattern: gains at mild budgets, losses where pruning pressure is
highest. Instrumentation explains the pattern. The coverage signal is
computed from a fixed sample of 8 training views and max-normalized; the
resulting per-anchor vector is extremely skewed (mean $\approx 0.000$: a
handful of anchors hold the entire normalized mass, the silent majority is
exactly zero). The additive score $0.7\,\tilde{c} + 0.3\,\tilde{s} +
0.25\,(a+u)$ therefore pins its scale on a few anchors and demotes the
calibrated ADMM score to a third-order term for the unseen majority---exactly
where ranking noise is executed most often, i.e., at the tightest budget.
Two structural guards remained effective throughout: the sensitivity cap
($w \le 0.3$) and the per-cell coverage constraint prevent the catastrophic
failure mode of an uncapped fusion (a $1.56$~dB collapse at
$\lambda{=}0.004$), bounding the damage to parity. Based on this analysis
the score was revised to a rotating-view coverage EMA (1080 distinct view
samples over training), percentile normalization, and a multiplicative
modulation $(a+u)\,(1+\gamma\,\tilde{\imath})$ that leaves the calibrated
ADMM ranking of zero-coverage anchors untouched. The revision thickens the
coverage distribution as designed (coverage $p50$ rises from $0.000$ to
$0.176$) and improves the tightest-budget point ($+0.06$~dB at
$\lambda{=}0.004$, now size-neutral), but does not clear our pre-registered
improvement gates ($\ge 27.46$ at $\lambda{=}0.004$,
$\ge 27.80$ at $\lambda{=}0.002$ on 1-78; the latter passed by $0.009$~dB).
We therefore position coverage-aware fusion pruning as a robustness
mechanism: its guards eliminate the collapse failure mode at an RD cost
that is within noise, rather than as a rate--distortion improvement.

\subsection{Negative results}
We document the directions that did not survive measurement. (i)~A DINO
semantic prior is predictable from decoder-recomputable features ($r=0.60$
with rich context) but is the wrong target: semantic identity does not
correspond to quantization importance, and the predictor gate initially
passed on an input set richer than the one the supervised network actually
receives. (ii)~An opacity-weighted reconstruction loss (per-pixel importance
from rendered alpha) degrades PSNR at two of three $\lambda$ settings. (iii)~MLP
weight quantization saves only $\sim$0.5\% of the payload and its original
implementation mixed quantized-size accounting with a float32 model; we
removed it rather than report the inconsistent number. (iv)~A sensitivity
side-information variant increases the payload with no quality gain and is
disabled. (v)~Hierarchical context adds size for negligible gain. These
negatives shaped the two retained mechanisms: every signal that survives is
either decoder-reproducible (quantization) or training-only (pruning).

\section{Conclusion}
\DCCA{} makes two decisions inside anchor-based 3DGS compression content-aware
without spending a bit of side information. A decoder-reproducible complexity
multiplier, supervised in training by render sensitivity, allocates quantization
precision where the image needs it. Coverage-aware fusion pruning, built on ADMM
budget sparsification, decides which anchors survive by what they cover---not only
by how complex they look---and removes the catastrophic low-rate failure mode at a
quality cost within measurement noise. Across three large-scale UAV scenes the
framework holds $58\%$ of the HAC++ model size at matched $\lambda$ within
$0.27$--$0.61$~dB, and aggregates the highest per-scene rank among 18 compared
methods. Both mechanisms respect the same
principle: every signal that shapes the bitstream must be recomputable from the
bitstream itself.

\bibliographystyle{IEEEtran}
\bibliography{references}

\end{document}
